\section{The Primal State}

\begin{quotex}
Race, caste, etc. e.g., gender, intelligence, strength (tr) exist in the spirit before manifesting in earthly and historical existence. Diversity has its origin from above, and that which is related to it on earth is only its reflection and symbol. As we willed to be on the basis of a primordial nature or a transcendental decision, so it is. It is not birth that determines character, but vice versa. It is character in its widest sense, because here common words are deceptive, that determines birth.
\flright{\textsc{Julius Evola}, \emph{Sintesi}}
\end{quotex}


The third visionary of the \emph{Examined Life} is a woman shown walking along a river in a chilly climate. She never tells use where she is going, although she knows how she got there: over a bridge that, as she points out, now has handicap access. While it is a good to provide aid to the disabled, since it is an accidental quality, the next step in her argument is more dubious. She compares building a wheelchair ramp, which is intended to enable normal activities, to the necessity to provide special advantages for women and minorities, the intent of which is the contrary.

Ignoring the best minds of the entire human race, she begins the discussion with ``Social Contract'' theory beginning with Locke, Kant, and Rousseau. This postulates a mythical past when individuals banded together to agree on a specific social arrangement, a ``contract'', to protect their individual interests.

This is contrary to all tradition for which the individual, as such, does not exist. A baby is born into a family, a social group, a political arrangement, and initiated into a spiritual tradition. Furthermore, that is all lawful as it was founded by a semi-divine figure who set down the law for it. Besides his natural well-being, the individual has a supernatural origin and destiny which is also supported by his family, tribe, and state.

So back to the professor. Given the originary formulation of the social contract, how then to account to the allegedly inferior position of women and minorities? Unable and unwilling to deny the obvious, she points out that individuals do not all have the same capability: some are simply more intelligent or stronger, hence they stack the deck in their favor. There is a German saying, ``Wer A sagt, muss auch B sagen,'' except it doesn't apply to the professor. She identifies the weak, vulnerable, or stupid, but leaves unsaid who the strong and intelligent are.

This impinges on the very notion of justice, whether the Fates are just or unjust. The modern view of justice is that everyone should be treated equally; the traditional view is that everyone gets what he deserves. Hence, in the ancient worlds, our place in life is not unjust. The group is arranged hierarchically, held together by organic bonds of love, loyalty, and fidelity. The child does not negotiate with his father for his benefits; rather the father provides for the child out of love and sense of duty. The stupid are not used by the intelligent; rather, from the knowledgeable, the dull are given fire, the wheel, the plough; they are informed when to plant and when to harvest. The weak are not opposed by the strong. The strong, in the traditional sense, are those who fight the greater battle; as initiates, they are knowers of the supernatural world. The weak cannot fight that inner battle on their own. They are given an exoteric faith, a set of rites, that provide the structure for them not to succumb to the forces of Chaos.

Yet, into this structure, the forces of subversion intrude, feeding themselves on fomenting envy, dissatisfaction, division, and resentment: the elements of Strife. It arises as an anti-intelligence, seemingly intelligent but more akin to slyness or cunning, that uses lies masquerading as truth as its lever. It is this subversion, in fact, that takes advantage of the vulnerable, weak, and stupid, while promising them that false allure of liberation.

So if, in the state of nature she imagines, the intelligent and the strong do not have the natural right to rule, or, better said, the duty to rule, then how can the weak and stupid negotiate their own share of power? Only through numerical superiority and intellectual subversion. The first step is to attack the supernatural foundations of the society. Those whose understanding is based only on belief are vulnerable to this approach. The second step is to overturn the organic bonds of love, loyalty, and fidelity, whose support has now been undermined, in order to create an unstructured mass of deracinated individuals. The third step is to unleash the forces of chaos which are superficially attractive and come to dominate the minds of the deracinated. Sensuality, materiality, novelty, the formless, the abnormal, the ugly, and the unleashing of instinct are celebrated in art, music and film.

So, in the professor's social contact, intelligence is replaced by cunning and individual strength guided by Order is replaced by the power of numbers dominated by Chaos. Rule by the Wise and Powerful is lost. As a matter of fact, the sophisticated do not even believe in such qualities. In a nutshell, instead of a Primordial State resulting in a fall from supernatural harmony, she envisions a Primal State of all against all, striving for personal gain, that will progressively evolve into some future utopia.


\flrightit{Posted on 2011-09-10 by Cologero}

\begin{center}* * *\end{center}

\begin{footnotesize}
\begin{sffamily}
\texttt{logres on 2011-09-10 at 16:18 said:}

Martha Nussbaum is a great example of the ``noble'' kind of modern philosopher (who incidentally condemned the film, but nevertheless subscribes to its postulates, or at least, the postulates it entails).

\hfill

\texttt{Matt on 2011-09-10 at 18:39 said:}

I think the only thinker in Examined Life who had something genuinely interesting to say was Zizek, in how everything that is nature isn't all wonderful and good as new agers and the modern environmentalist movement would like people to think. This starting point is in the spirit of traditional thinking, however, tradition teaches reality extends beyond mere nature and there is that which allows the individual to experience stability (grace and personality, for example). Zizek, as a materialist, can't go beyond that starting point, and so one is left with an inadequate (possibly counterfeit) way of dealing with this fact of the natural world.


\hfill

\end{sffamily}
\end{footnotesize}

